%%═══════════════════════════════════════════════════════════════════
%% Lecture 03 — Canonical Quantisation of the Real Scalar Field:
%%              Conjugate Momentum, Equal-Time Commutators,
%%              Fourier Decomposition, and Creation/Annihilation Operators
%% 77609 — Introduction to Elementary Particles
%% Date: 12/04/2026
%% Lecturer: Prof. Yonit Hochberg
%%═══════════════════════════════════════════════════════════════════

\section{Lecture 03 — Canonical Quantisation of the Real Scalar Field}\label{sec:lec03:quantisation}
\begin{center}
\begin{tabular}{ll}
  \textbf{Date:}\quad 12/04/2026   & \\
\end{tabular}
\end{center}
\hrule\vspace{1.5em}
%%──────────────────────────────────────────────────────────────────
\subsection*{Overview}

Having established classical field theory in Lecture~02, we now
perform the pivotal step of \textbf{canonical quantisation}: promoting
the classical real scalar field $\phi(x,t)$ and its conjugate momentum
density $\pi(x,t)$ to \emph{operators} satisfying the
\textbf{equal-time canonical commutation relations} (ET-CCR).

%% Key boxed result: equal-time CCR
\begin{equation}
  \boxed{%
    \bigl[\phi(\mathbf{x},t),\,\pi(\mathbf{y},t)\bigr]
    = i\,\delta^{(3)}(\mathbf{x}-\mathbf{y}),
    \qquad
    \bigl[\phi(\mathbf{x},t),\,\phi(\mathbf{y},t)\bigr]
    = \bigl[\pi(\mathbf{x},t),\,\pi(\mathbf{y},t)\bigr]=0.
  }
  \label{eq:lec03:ETCCR}
\end{equation}

The Klein-Gordon equation, linear in $\phi$, calls for a
Fourier decomposition in momentum space.  For each Fourier mode
$\mathbf{p}$, the field satisfies a simple harmonic oscillator equation
with frequency $\omega_{\mathbf{p}}=\sqrt{|\mathbf{p}|^2+m^2}$
(the relativistic dispersion relation).  The reality condition on the
classical field, and the hermiticity condition on the quantum field,
constrain the Fourier coefficients: they become the familiar creation
and annihilation operators $a_{\mathbf{p}}^\dagger$ and
$a_{\mathbf{p}}$.  This lecture thereby establishes the bridge between
the classical field and the Fock-space picture of free particles.

%%──────────────────────────────────────────────────────────────────
\subsection*{Step 1: Conjugate momentum density from the Lagrangian}

We work in natural units ($\hbar=c=1$) and in $3+1$ spacetime
dimensions.  The free real scalar field Lagrangian density is
\begin{equation}
  \lagr
  = \tfrac{1}{2}\,\partial_\mu\phi\,\partial^\mu\phi - \tfrac{1}{2}m^2\phi^2
  = \tfrac{1}{2}\bigl(\dphi^2 - |\nabla\phi|^2\bigr) - \tfrac{1}{2}m^2\phi^2.
  \label{eq:lec03:L}
\end{equation}

\dfn{Conjugate momentum density}{By analogy with analytical mechanics—where the conjugate momentum to
$q$ is $p=\partial L/\partial\dot{q}$—the conjugate momentum density
associated with the field $\phi$ is
\begin{equation}
  \pi(x,t)
  \equiv \frac{\partial\lagr}{\partial(\partial_t\phi)}
  = \dot\phi(x,t).
  \label{eq:lec03:pi}
\end{equation}}

\begin{remark}[Density vs.\ total momentum]
The quantity $\pi$ defined above is a \emph{density} (it derives from
the Lagrangian density $\lagr$, not the integrated Lagrangian $L$).
This is the standard convention in QFT.  The momentum $p$ in the
Fourier decomposition below is the \emph{ordinary} (kinematic) 3-momentum
label; it should not be confused with $\pi$.
\end{remark}

%%──────────────────────────────────────────────────────────────────
\subsection*{Step 2: Hamiltonian density}

The Hamiltonian density (via the Legendre transform of $\lagr$) is

%% Derivation: Hamiltonian density from Legendre transform
\begin{align}
  \hamilt
  &= \pi\,\dphi - \lagr \notag \\
  &= \dot\phi^2
     - \Bigl[\tfrac{1}{2}\dot\phi^2
             - \tfrac{1}{2}|\nabla\phi|^2
             - \tfrac{1}{2}m^2\phi^2\Bigr] \notag \\
  &= \tfrac{1}{2}\dot\phi^2
     + \tfrac{1}{2}|\nabla\phi|^2
     + \tfrac{1}{2}m^2\phi^2.
  \label{eq:lec03:H}
\end{align}

Using $\pi=\dphi$ we write this in the canonical form
\begin{equation}
  \hamilt
  = \tfrac{1}{2}\pi^2
    + \tfrac{1}{2}|\nabla\phi|^2
    + \tfrac{1}{2}m^2\phi^2,
\end{equation}
and the total Hamiltonian is
$H = \int d^3x\;\hamilt(\mathbf{x},t)$.

%%──────────────────────────────────────────────────────────────────
\subsection*{Step 3: The quantisation prescription}

\dfn{Canonical quantisation of fields}{We \emph{promote} $\phi$ and $\pi$ from classical functions to
Hermitian operators on a Hilbert space and impose the
\textbf{equal-time canonical commutation relations}:
\begin{align}
  \bigl[\phi(\mathbf{x},t),\;\pi(\mathbf{y},t)\bigr]
    &= i\,\delta^{(3)}(\mathbf{x}-\mathbf{y}),
    \label{eq:lec03:CCR-phi-pi}\\
  \bigl[\phi(\mathbf{x},t),\;\phi(\mathbf{y},t)\bigr]
    &= 0,
    \label{eq:lec03:CCR-phi-phi}\\
  \bigl[\pi(\mathbf{x},t),\;\pi(\mathbf{y},t)\bigr]
    &= 0.
    \label{eq:lec03:CCR-pi-pi}
\end{align}
These are the direct field-theory analogues of $[q_i,p_j]=i\delta_{ij}$.}

\begin{remark}[Why equal-time?]
The commutators are imposed at the \emph{same time} $t$.  This singles
out a time slice (a Cauchy surface) and is the natural generalization of
the single-time quantum mechanics condition $[q(t),p(t)]=i\hbar$.
Although equal-time quantisation breaks manifest Lorentz covariance, the
resulting QFT is nonetheless fully Lorentz invariant once interactions
are included properly.  A Lorentz-covariant quantisation
exists but is more involved.
\end{remark}

%%──────────────────────────────────────────────────────────────────
\subsection*{Step 4: Fourier decomposition and the Klein-Gordon spectrum}

%% Derivation: Fourier decomposition of the scalar field
The Klein-Gordon equation arising from \eqref{eq:lec03:L} via
the Euler-Lagrange prescription is
\begin{equation}
  (\Box + m^2)\phi = 0,
  \qquad
  \Box \equiv \partial^\mu\partial_\mu
             = \partial_t^2 - \nabla^2.
  \label{eq:lec03:KG}
\end{equation}

\subsubsection*{Fourier decomposition in momentum space}

Write the field as a Fourier integral over 3-momentum:
\begin{equation}
  \phi(\mathbf{x},t)
  = \int \frac{d^3p}{(2\pi)^3}\;
       \tilde\phi(\mathbf{p},t)\,e^{i\mathbf{p}\cdot\mathbf{x}}.
  \label{eq:lec03:Fourier-def}
\end{equation}
Substituting into \eqref{eq:lec03:KG}:
\begin{align}
  0
  &= (\partial_t^2 - \nabla^2 + m^2)\phi \notag\\
  &= \int \frac{d^3p}{(2\pi)^3}
       \bigl[\,\ddot{\tilde\phi}(\mathbf{p},t)
             + |\mathbf{p}|^2\,\tilde\phi(\mathbf{p},t)
             + m^2\,\tilde\phi(\mathbf{p},t)\bigr]
       e^{i\mathbf{p}\cdot\mathbf{x}}.
  \label{eq:lec03:KG-decomp}
\end{align}
Since this must hold for every mode independently,
\begin{equation}
  \ddot{\tilde\phi}(\mathbf{p},t)
  + \underbrace{\bigl(|\mathbf{p}|^2 + m^2\bigr)}_{\omega_{\mathbf{p}}^2}
    \tilde\phi(\mathbf{p},t) = 0.
  \label{eq:lec03:SHO}
\end{equation}

\dfn{Relativistic dispersion relation}{Define the \emph{on-shell frequency} for each mode:
\begin{equation}
  \omega_{\mathbf{p}}
  \equiv \sqrt{|\mathbf{p}|^2 + m^2}.
  \label{eq:lec03:omega}
\end{equation}
This is precisely the relativistic energy-momentum relation $E^2=p^2+m^2$
(in natural units) for a particle of mass $m$ and 3-momentum $\mathbf{p}$.}

Equation~\eqref{eq:lec03:SHO} is the equation of a
\textbf{simple harmonic oscillator} with frequency $\omega_{\mathbf{p}}$
for each Fourier mode.  The general solution is
\begin{equation}
  \tilde\phi(\mathbf{p},t)
  = A(\mathbf{p})\,e^{-i\omega_{\mathbf{p}}t}
    + B(\mathbf{p})\,e^{+i\omega_{\mathbf{p}}t}.
  \label{eq:lec03:SHO-sol}
\end{equation}

%%──────────────────────────────────────────────────────────────────
\subsection*{Step 5: Reality (hermiticity) condition on the field}

%% Derivation: Reality condition constrains Fourier coefficients
The \emph{classical} field $\phi(\mathbf{x},t)$ is real:
$\phi^*(\mathbf{x},t)=\phi(\mathbf{x},t)$.  Taking the complex
conjugate of the Fourier representation and requiring equality gives
\begin{align}
  \phi^*(\mathbf{x},t)
  &= \int \frac{d^3p}{(2\pi)^3}
        \tilde\phi^*(\mathbf{p},t)\,e^{-i\mathbf{p}\cdot\mathbf{x}}
  \notag\\
  &= \int \frac{d^3p}{(2\pi)^3}
        \tilde\phi^*(-\mathbf{p},t)\,e^{+i\mathbf{p}\cdot\mathbf{x}}
        \quad [\text{re-label }\mathbf{p}\to -\mathbf{p}],
  \label{eq:lec03:reality-step}
\end{align}
where in the last line we used the substitution $\mathbf{p}\to-\mathbf{p}$
to match the $e^{i\mathbf{p}\cdot\mathbf{x}}$ Fourier basis.
Comparing with \eqref{eq:lec03:Fourier-def} we need
\begin{equation}
  \tilde\phi(\mathbf{p},t) = \tilde\phi^*(-\mathbf{p},t).
  \label{eq:lec03:reality-cond}
\end{equation}
Applying this to \eqref{eq:lec03:SHO-sol}:
\begin{align}
  A(\mathbf{p})e^{-i\omega_{\mathbf{p}} t}
  + B(\mathbf{p})e^{+i\omega_{\mathbf{p}} t}
  &= A^*(-\mathbf{p})e^{+i\omega_{\mathbf{p}} t}
   + B^*(-\mathbf{p})e^{-i\omega_{\mathbf{p}} t}.
  \label{eq:lec03:reality-match}
\end{align}
Matching the $e^{-i\omega t}$ and $e^{+i\omega t}$ coefficients
separately:
\begin{equation}
  B(\mathbf{p}) = A^*(-\mathbf{p}).
  \label{eq:lec03:B-from-A}
\end{equation}
Hence the general real solution contains only one independent
complex amplitude per mode, which we rename $a(\mathbf{p})\equiv A(\mathbf{p})$:
\begin{equation}
  \tilde\phi(\mathbf{p},t)
  = a(\mathbf{p})\,e^{-i\omega_{\mathbf{p}} t}
    + a^*(-\mathbf{p})\,e^{+i\omega_{\mathbf{p}} t}.
  \label{eq:lec03:phi-tilde}
\end{equation}

%%──────────────────────────────────────────────────────────────────
\subsection*{Step 6: Full mode expansion of the quantum field}

Substituting \eqref{eq:lec03:phi-tilde} back into the Fourier integral
\eqref{eq:lec03:Fourier-def} and relabelling $\mathbf{p}\to -\mathbf{p}$
in the second term:
\begin{align}
  \phi(\mathbf{x},t)
  &= \int \frac{d^3p}{(2\pi)^3}
     \Bigl[
       a(\mathbf{p})\,e^{-i\omega_{\mathbf{p}} t + i\mathbf{p}\cdot\mathbf{x}}
       + a^*(-\mathbf{p})\,e^{+i\omega_{\mathbf{p}} t + i\mathbf{p}\cdot\mathbf{x}}
     \Bigr]
  \notag \\
  &= \int \frac{d^3p}{(2\pi)^3}
     \Bigl[
       a(\mathbf{p})\,e^{-i\omega_{\mathbf{p}} t + i\mathbf{p}\cdot\mathbf{x}}
       + a^*(\mathbf{p})\,e^{+i\omega_{\mathbf{p}} t - i\mathbf{p}\cdot\mathbf{x}}
     \Bigr].
  \label{eq:lec03:phi-expand-pre}
\end{align}
It is conventional to include a Lorentz-invariant normalisation factor
$1/(2\omega_{\mathbf{p}})^{1/2}$, giving the standard mode expansion:
\begin{equation}
  \boxed{
  \phi(\mathbf{x},t)
  = \int \frac{d^3p}{(2\pi)^3}
    \frac{1}{\sqrt{2\omega_{\mathbf{p}}}}
    \Bigl[
      a_{\mathbf{p}}\,e^{-i(\omega_{\mathbf{p}} t - \mathbf{p}\cdot\mathbf{x})}
      + a_{\mathbf{p}}^\dagger\,
        e^{+i(\omega_{\mathbf{p}} t - \mathbf{p}\cdot\mathbf{x})}
    \Bigr].
  }
  \label{eq:lec03:phi-mode}
\end{equation}
In the quantum theory $a(\mathbf{p})\to a_{\mathbf{p}}$ and
$a^*(\mathbf{p})\to a_{\mathbf{p}}^\dagger$ become
\textbf{operators} on the Fock space;
the reality condition $\phi=\phi^*$ becomes the hermiticity condition
$\phi=\phi^\dagger$.

Similarly, the conjugate momentum $\pi=\dphi$ has the expansion
\begin{equation}
  \pi(\mathbf{x},t)
  = \int \frac{d^3p}{(2\pi)^3}
    \frac{(-i)\sqrt{\omega_{\mathbf{p}}}}{\sqrt{2}}
    \Bigl[
      a_{\mathbf{p}}\,e^{-i(\omega_{\mathbf{p}} t - \mathbf{p}\cdot\mathbf{x})}
      - a_{\mathbf{p}}^\dagger\,
        e^{+i(\omega_{\mathbf{p}} t - \mathbf{p}\cdot\mathbf{x})}
    \Bigr],
  \label{eq:lec03:pi-mode}
\end{equation}
obtained by differentiating \eqref{eq:lec03:phi-mode} with respect to $t$.

%%──────────────────────────────────────────────────────────────────
\subsection*{Step 7: Commutation relations for $a_{\mathbf{p}}$ and $a_{\mathbf{p}}^\dagger$}

%% Derivation: [a,a†] from ET-CCR
Imposing the ET-CCR \eqref{eq:lec03:ETCCR} on the mode expansions
\eqref{eq:lec03:phi-mode}--\eqref{eq:lec03:pi-mode}, one finds (by
inverting the Fourier transform and using the completeness of plane waves)
that the creation and annihilation operators satisfy
\begin{equation}
  \bigl[a_{\mathbf{p}},\,a_{\mathbf{q}}^\dagger\bigr]
  = (2\pi)^3\,\delta^{(3)}(\mathbf{p}-\mathbf{q}),
  \label{eq:lec03:a-comm}
\end{equation}
\begin{equation}
  \bigl[a_{\mathbf{p}},\,a_{\mathbf{q}}\bigr]
  = \bigl[a_{\mathbf{p}}^\dagger,\,a_{\mathbf{q}}^\dagger\bigr]
  = 0.
  \label{eq:lec03:a-zero-comm}
\end{equation}
These are precisely the commutation relations of an
\textbf{infinite collection of decoupled quantum harmonic oscillators},
one for each momentum mode $\mathbf{p}$.

\begin{remark}[Physical interpretation of $a_{\mathbf{p}}$]
Before quantisation, $a(\mathbf{p})$ is merely a complex Fourier
coefficient—a \emph{c-number}.  After quantisation it becomes an
\emph{operator}: $a_{\mathbf{p}}$ \textbf{annihilates} a particle of
momentum $\mathbf{p}$, while $a_{\mathbf{p}}^\dagger$
\textbf{creates} one.  The Fock vacuum $|0\rangle$ is defined by
$a_{\mathbf{p}}|0\rangle=0$ for all $\mathbf{p}$.
\end{remark}

%%──────────────────────────────────────────────────────────────────
\subsection*{Examples}

\ex{Explicit derivation of $\pi$ from $\phi$}{Starting from the mode expansion of $\phi$ in~\eqref{eq:lec03:phi-mode},
differentiate with respect to $t$:
\begin{align}
  \pi(\mathbf{x},t) = \dphi
  &= \int \frac{d^3p}{(2\pi)^3}
     \frac{1}{\sqrt{2\omega_{\mathbf{p}}}}
     \frac{\partial}{\partial t}
     \Bigl[
       a_{\mathbf{p}}\,e^{-i(\omega_{\mathbf{p}} t - \mathbf{p}\cdot\mathbf{x})}
       + a_{\mathbf{p}}^\dagger\,e^{+i(\omega_{\mathbf{p}} t - \mathbf{p}\cdot\mathbf{x})}
     \Bigr] \notag\\
  &= \int \frac{d^3p}{(2\pi)^3}
     \frac{1}{\sqrt{2\omega_{\mathbf{p}}}}
     \Bigl[
       (-i\omega_{\mathbf{p}})\,a_{\mathbf{p}}\,e^{-i(\omega_{\mathbf{p}} t - \mathbf{p}\cdot\mathbf{x})}
       + (i\omega_{\mathbf{p}})\,a_{\mathbf{p}}^\dagger\,e^{+i(\omega_{\mathbf{p}} t - \mathbf{p}\cdot\mathbf{x})}
     \Bigr] \notag\\
  &= -i\int \frac{d^3p}{(2\pi)^3}
     \sqrt{\frac{\omega_{\mathbf{p}}}{2}}
     \Bigl[
       a_{\mathbf{p}}\,e^{-i(\omega_{\mathbf{p}} t - \mathbf{p}\cdot\mathbf{x})}
       - a_{\mathbf{p}}^\dagger\,e^{+i(\omega_{\mathbf{p}} t - \mathbf{p}\cdot\mathbf{x})}
     \Bigr].
\end{align}
This confirms the expression in~\eqref{eq:lec03:pi-mode}.}

\ex{Each mode as a harmonic oscillator}{Fix a mode $\mathbf{p}$ and define the ``position'' and ``momentum'' of
that mode:
\begin{equation}
  Q_{\mathbf{p}} \propto a_{\mathbf{p}} + a_{\mathbf{p}}^\dagger,
  \qquad
  P_{\mathbf{p}} \propto -i(a_{\mathbf{p}} - a_{\mathbf{p}}^\dagger).
  \label{eq:lec03:mode-QP}
\end{equation}
From \eqref{eq:lec03:a-comm} one recovers $[Q_{\mathbf{p}},P_{\mathbf{p}}]=i$,
and the Hamiltonian contribution of this mode is

\begin{equation}
  H_{\mathbf{p}}
  = \omega_{\mathbf{p}}
    \Bigl(a_{\mathbf{p}}^\dagger a_{\mathbf{p}} + \tfrac{1}{2}\Bigr).
  \label{eq:lec03:mode-H}
\end{equation}

This is identical to a quantum harmonic oscillator of frequency
$\omega_{\mathbf{p}}=\sqrt{|\mathbf{p}|^2+m^2}$.  The total Hamiltonian is
therefore
\begin{equation}
  H = \int \frac{d^3p}{(2\pi)^3}\;
      \omega_{\mathbf{p}}
      \Bigl(a_{\mathbf{p}}^\dagger a_{\mathbf{p}} + \tfrac{1}{2}\Bigr),
  \label{eq:lec03:H-normal}
\end{equation}
where the $\tfrac{1}{2}$ term gives the (formally infinite) zero-point
energy of the vacuum, which is later handled by normal ordering.}

\ex{Dispersion relation as the relativistic energy-momentum relation}{The on-shell frequency defined in \eqref{eq:lec03:omega},
$\omega_{\mathbf{p}} = \sqrt{|\mathbf{p}|^2 + m^2}$,
is equivalent to $E^2 = p^2 + m^2$ (natural units).
For a massless field ($m=0$) this reduces to $\omega_{\mathbf{p}}=|\mathbf{p}|$
(linear dispersion, i.e.\ photons or massless scalar).
For a non-zero mass it implies a \textbf{minimum energy} $E_{\min}=m$
at $\mathbf{p}=0$, corresponding to a particle at rest.}

%%──────────────────────────────────────────────────────────────────
%% [BOARD CONTENT NOT CAPTURED — Lec03 timestamp 00:04:13–00:08:49]
%% (Classroom logistics, seating, and transition to physics content)

%%──────────────────────────────────────────────────────────────────
\subsection*{To Remember}

\begin{itemize}
  \item The conjugate momentum density of the real scalar field is
    $\pi = \partial\lagr/\partial\dphi = \dphi$.
  \item Canonical quantisation promotes $\phi$ and $\pi$ to Hermitian
    operators satisfying the ET-CCR:
    $[\phi(\mathbf{x},t),\pi(\mathbf{y},t)]=i\delta^{(3)}(\mathbf{x}-\mathbf{y})$.
  \item The Klein-Gordon equation in each Fourier mode reduces to a
    simple harmonic oscillator equation with frequency
    $\omega_{\mathbf{p}}=\sqrt{|\mathbf{p}|^2+m^2}$ (relativistic dispersion).
  \item The reality (hermiticity) condition on $\phi$ constrains the two
    Fourier coefficients to be complex conjugates:
    $B(\mathbf{p})=A^*(-\mathbf{p})$, leaving one independent operator
    per mode.
  \item The mode expansion is
    $\phi = \int \frac{d^3p}{(2\pi)^3}\frac{1}{\sqrt{2\omega_{\mathbf{p}}}}
    \bigl(a_{\mathbf{p}}\,e^{-ipx}+a_{\mathbf{p}}^\dagger e^{+ipx}\bigr)$,
    where $px\equiv\omega_{\mathbf{p}}t - \mathbf{p}\cdot\mathbf{x}$.
  \item Imposing the ET-CCR on the mode expansion yields the
    harmonic-oscillator algebra:
    $[a_{\mathbf{p}}, a_{\mathbf{q}}^\dagger]=(2\pi)^3\delta^{(3)}(\mathbf{p}-\mathbf{q})$.
  \item The full Hamiltonian is a sum (integral) of harmonic oscillators:
    $H=\int\frac{d^3p}{(2\pi)^3}\,\omega_{\mathbf{p}}
    \bigl(a_{\mathbf{p}}^\dagger a_{\mathbf{p}}+\tfrac{1}{2}\bigr)$;
    the infinite zero-point energy is removed by normal ordering.
\end{itemize}

%%──────────────────────────────────────────────────────────────────
\subsection*{Further Reading}

No specific textbook was cited in this lecture.  The following cover the same
material at varying depths:

\begin{itemize}
  \item \textbf{Peskin \& Schroeder, \S2.1--2.3} — The canonical quantization
    of the real Klein--Gordon field from first principles: mode expansion,
    creation/annihilation operators, Fock space, normal ordering, and the Feynman
    propagator.  This is the primary QFT reference for the material in this
    lecture.

  \item \textbf{Tong, \textit{QFT Notes}, Chapter~2}
    (\href{https://www.damtp.cam.ac.uk/user/tong/qft.html}{\texttt{damtp.cam.ac.uk/user/tong/qft.html}}):
    ``Free Fields'' — covers the same quantization procedure in a more pedagogical
    style, with explicit derivations of the commutation relations and the Hamiltonian
    as a sum of harmonic oscillators.

  \item \textbf{Griffiths, \textit{Introduction to Elementary Particles}, Appendix}
    — Brief summary of creation and annihilation operators and their algebra, useful
    for a quick reference on the operator formalism.

  \item \textbf{Lancaster \& Blundell, \textit{Quantum Field Theory for the Gifted
    Amateur}, Chapters~11--12}
    (\href{https://global.oup.com/academic/product/quantum-field-theory-for-the-gifted-amateur-9780199699322}{\texttt{OUP}}):
    a very readable step-by-step quantization of the scalar field,
    explicitly connecting each step to the quantum harmonic oscillator the student
    already knows.
\end{itemize}
